\documentclass[11pt]{article}

% ── Encoding & fonts ──────────────────────────────────────────────────────────
\usepackage[T1]{fontenc}
\usepackage[utf8]{inputenc}
\usepackage{lmodern}
\usepackage{microtype}

% ── Page geometry ─────────────────────────────────────────────────────────────
\usepackage[margin=1.25in]{geometry}

% ── Mathematics ───────────────────────────────────────────────────────────────
\usepackage{amsmath}
\usepackage{amssymb}

% ── Tables ────────────────────────────────────────────────────────────────────
\usepackage{booktabs}
\usepackage{array}
\usepackage{tabularx}

% ── Lists ─────────────────────────────────────────────────────────────────────
\usepackage{enumitem}
\setlist{itemsep=2pt, topsep=4pt}

% ── Code / typewriter text ────────────────────────────────────────────────────
\usepackage{listings}
\lstset{basicstyle=\ttfamily\small, breaklines=true}

% ── Hyperlinks ────────────────────────────────────────────────────────────────
\usepackage[colorlinks=true, linkcolor=black, citecolor=black,
            urlcolor=black, pdfborder={0 0 0}]{hyperref}
\usepackage{url}

% ── Bibliography ──────────────────────────────────────────────────────────────
\usepackage[numbers, sort&compress]{natbib}
\bibliographystyle{unsrtnat}

% ── Paragraph spacing ─────────────────────────────────────────────────────────
\usepackage{parskip}
\setlength{\parskip}{6pt}

% ── Section formatting ────────────────────────────────────────────────────────
\usepackage{titlesec}
\titleformat{\section}{\large\bfseries}{\thesection.}{0.6em}{}
\titleformat{\subsection}{\normalsize\bfseries}{\thesubsection}{0.6em}{}
\titlespacing{\section}{0pt}{14pt}{6pt}
\titlespacing{\subsection}{0pt}{10pt}{4pt}

% ── Abstract ──────────────────────────────────────────────────────────────────
\usepackage{abstract}
\renewcommand{\abstractnamefont}{\normalfont\bfseries}
\renewcommand{\abstracttextfont}{\normalfont\small}

% ── Footnotes ─────────────────────────────────────────────────────────────────
\usepackage[hang, flushmargin]{footmisc}

% ─────────────────────────────────────────────────────────────────────────────
\title{\textbf{The Coordinator Trap: Structural Scaling Liabilities in\\
Mediated Multi-Agent Architectures and a Substrate-Based Alternative}}

\author{%
  Dr.\ Richard Nicholson\\[4pt]
  \small Tathata Systems Ltd, United Kingdom\\
  \small \href{mailto:tathatasystems@proton.me}{tathatasystems@proton.me}
}

\date{June 2026\\[4pt]
\small Preprint. \copyright\ 2026 Tathata Systems Ltd. Licensed under
\href{https://creativecommons.org/licenses/by/4.0/}{CC BY 4.0}.\\
\small DOI: \href{https://doi.org/10.5281/zenodo.20665238}{10.5281/zenodo.20665238}}
% ─────────────────────────────────────────────────────────────────────────────

\begin{document}

\maketitle

\begin{abstract}
Multi-agent coordination systems are converging on two dominant architectural
patterns: \emph{mediated hierarchies}, in which a central engine aggregates agent
outputs, decides, and issues commands; and \emph{registry-based discovery}, in
which a global index brokers introductions between agents. Both patterns introduce
a coordinator---by different names---that becomes a bottleneck, a single point of
failure, and an unbounded audit obligation as agent populations grow. We
argue that these failure modes are not implementation deficiencies
but structural consequences of the coordinator assumption itself. Drawing on
Holland's framework for Complex Adaptive Systems, we propose an alternative in
which coordination emerges from substrate properties rather than explicit
protocols. We present Mycelium, a working implementation of this model as an
embeddable Rust library, and show how its three-layer architecture---gossip KV
store, signal mesh with boundary admission, and epidemic consensus---structurally
addresses each failure mode rather than ameliorating it.
\end{abstract}

% ─────────────────────────────────────────────────────────────────────────────
\section{Introduction}

A dominant family of AI agent frameworks released between 2023 and 2026---the
supervisor, manager, hierarchical orchestration, and router patterns found in
LangChain, AutoGen, CrewAI, and their successors---converges on the same
structural pattern: a central orchestrator receives requests, routes them to
specialist sub-agents, aggregates their outputs, and returns a synthesised
result~\cite{silobench,mas-survey,raps}. The orchestrator is the coordinator.
The sub-agents are its workers.

This paper argues that the coordinator is not a solution to the multi-agent
coordination problem. It is a restatement of the problem at a different scale.

The coordinator trap produces three failure modes that are structural, not
incidental: (i)~unbounded audit obligation, as every agent output passes through
the coordinator and generates an artifact of equivalent apparent authority
regardless of relevance; (ii)~context loss, as shared state held in or near the
coordinator is lost on restart, forcing agents to re-litigate settled decisions;
and (iii)~output format mismatch, as the coordinator broadcasts undifferentiated
results to heterogeneous consumers it has no model of.

These failure modes have been documented in production deployments. In the
Cisco Mycelium project---an unrelated multi-agent orchestration system---the
coordinator could not maintain context across restarts, produced three-hour
and five-second problems with the same apparent authority, and generated
structured machine logs that failed to communicate to human
reviewers~\cite{signal-noise,context-problem}.

We show that these failure modes are irreducible by improving the coordinator.
They follow necessarily from its existence. We then derive a coordinator-free
architecture from John Holland's signal/boundary model~\cite{holland} and present
Mycelium, a working implementation as an embeddable Rust library.

% ─────────────────────────────────────────────────────────────────────────────
\section{Background: The Two Dominant Patterns}

\subsection{Mediated Hierarchy}

A mediated hierarchy places a cognitive engine at the centre of the agent
population. Specialist agents receive structured payloads from the engine, process
them, and return results. The engine aggregates, decides, and broadcasts. Agent
populations in contemporary LLM frameworks~\cite{mas-survey} follow this pattern:
the orchestrator holds the conversation state and the routing logic; sub-agents
are workers in a fanout RPC system.

Kubernetes~\cite{kubernetes} represents a more mature instance of the same
architectural pattern. A control plane holds the declared target state;
reconciliation loops continuously compare actual cluster state to that target and
issue corrective actions. The coordinator is more sophisticated than a negotiation
engine, but the structural relationship is identical: state is held above the
runtime, and correctness is enforced top-down.

Independent empirical work corroborates that this pattern's failure is structural
rather than incidental. Rodriguez~\cite{pressure-fields-decay}, comparing
coordination strategies for LLM agents on constraint-satisfaction tasks, reports
that hierarchical (mediated) control applied \emph{zero} accepted changes in
$66.7\%$ of runs---a $98.7\%$ rejection rate---as the coordinator repeatedly
re-selected the single highest-pressure sub-problem while progress stalled
everywhere else. Notably, that study reaches for the same alternative this paper
argues for: it proposes a stigmergic, coordinator-free coordination scheme, and
finds it outperforms mediated control by roughly $30\times$ on the same task. It
does so, however, as an \emph{application-level} algorithm over a single shared
artifact on one machine---with a central per-tick selection queue, potential-game
(not distributed) convergence, and (its own §7.6) multi-artifact and larger-scale
operation left to future work. This paper supplies what that direction assumes but
does not build: the distributed substrate---causal state, receiver-side signal
boundaries, epidemic consensus, and recallable-role failover---developed in
Sections 5--7.

\subsection{Registry-Based Discovery}

A registry-based system separates coordination from discovery. Agents register
capabilities in a global index; other agents query the index to locate peers. The
registry brokers introductions; agents communicate point-to-point thereafter.
FIPA's Directory Facilitator~\cite{fipa,jade} is a canonical instance. MIT
NANDA~\cite{nanda} represents a more sophisticated point on the same spectrum:
it uses CRDT-based updates, cryptographically verifiable AgentFacts, and
multi-endpoint routing rather than a static index. This spectrum runs from
static centralised registries through federated and CRDT-verifiable designs to
fully local self-admission substrates. The claim here applies to the common
structural property they share: an authority-bearing component through which
discovery passes and which holds canonical state about the agent population.
Mycelium sits at the far end of this spectrum---no registry, no authority, no
discovery component---but the critique does not apply equally to every design
in the family.

\subsection{The Shared Assumption}

Superficially, these two patterns are different. A mediated hierarchy coordinates
behaviour at runtime; a registry brokers introductions at discovery time. But they
share a single structural assumption: \textbf{there exists a designated component
that holds authoritative state about the agent population and through which
coordination passes.}

This assumption is the coordinator trap. Section~\ref{sec:trap} traces each
failure mode to this assumption directly.

\subsection{The Separation Tax}

The coordinator assumption has a concrete operational cost. A representative
conventional ML inference pipeline must address at minimum four distinct
infrastructure concerns: task dispatch and queuing (Redis/Celery or Kafka
depending on scale), worker registration and discovery (Consul, etcd, or a
sidecar), failure recovery and dead-letter handling, and cross-component
observability. Each concern requires a separately deployed, operated, and
monitored service. The pipeline is not the model---it is the scaffolding around
the model.

Each concern is a separate operational surface: a separate failure mode, a
separate scaling policy, a separate upgrade cycle, and a conceptual translation
layer. The coordination work---route an item to a capable worker, throttle when
workers are saturated, clean up abandoned items, make topology visible---is the
same abstract work at every level. The \emph{separation tax} is the cost of
performing that work redundantly at each concern boundary rather than once in a
unified substrate.

The two dominant patterns described in \S\S2.1--2.2 do not reduce this tax. They
add another coordination system on top of the existing stack. This paper's claim
is that the separation tax is not a necessary cost of multi-agent coordination. It
is the cost of applying the coordinator assumption independently at each layer of
the stack.

% ─────────────────────────────────────────────────────────────────────────────
\section{A History of the Coordinator Assumption}

The coordinator assumption has been the default in distributed computing for over
fifty years. Each generation has recognised the problems it causes and attempted
to mitigate them without eliminating the assumption itself.

\subsection{Blackboard Systems (1970s--80s)}

The blackboard model, most prominently realised in HEARSAY-II~\cite{hearsay},
organises coordination around a central shared data structure and a set of
specialised knowledge sources that read from and write to it. A control component
monitors the blackboard and decides which knowledge source to activate next.
HEARSAY demonstrated that the model could solve hard problems; it also
demonstrated the failure modes: the control component became a bottleneck and a
single point of failure, and the blackboard accumulated stale entries.

\subsection{The Actor Model (Hewitt, 1973)}

Hewitt, Bishop and Steiger's Actor model~\cite{hewitt} provided the first formal
model of concurrent computation as communicating entities. The model eliminates
shared mutable state---a genuine advance---but retains explicit addressing: to
send a message, an actor must know the recipient's address. Topology must be
managed explicitly by the sender. The Actor model correctly identified agents as
the unit of composition; it did not identify boundaries as the mechanism for
self-organising communication.

Mature actor runtimes---most notably Erlang/OTP---address fault tolerance and
process supervision without a central coordinator, and represent a genuine step
beyond the original formulation. The addressing problem persists, however:
Erlang's process registry is a coordinator with a narrower mandate, and OTP
supervision trees impose a coordinator hierarchy for failure propagation.

\subsection{Linda and the Tuple Space (Gelernter, 1985)}

Gelernter's Linda~\cite{linda} introduced generative communication via a shared
tuple space. Carriero and Gelernter~\cite{linda-coord} established that
coordination is a distinct language design concern, separable from
computation---a correct and underappreciated insight. What Linda retained,
however, is a shared coordinator: the tuple space itself must remain accessible,
consistent, and available to all participants. Fault-tolerant variants required
replication schemes and consistency protocols. The medium is the coordinator. The
coordinator trap persists in a subtler form.

The distinction that matters---and that Section~7 will realise---is precise: in
Linda, the admission decision is a retrieval pattern issued against the shared
medium, which must remain consistently accessible to all participants. An
alternative is for each node to hold its own boundary as a local declaration
evaluated against live distributed state: the medium is never queried, the node
self-filters, and there is no shared state to keep consistent.

Two distinct things are entangled in the classical tuple-space critique, and
separating them recovers what Linda got right. The space is a
\emph{decision-maker} in no sense at all---it allocates nothing and predicts
nothing; workers claim work when ready, which is the one coordination act that
requires no knowledge of any other participant's state. The space is, however, a
\emph{rendezvous point on the data path}---a capacity and availability concern,
like any shared medium. The coordinator trap is a claim about the first role, not
the second; FT-Linda and JavaSpaces struggled because they answered the second
concern with designed-in consistency machinery. Mycelium's
\texttt{mycelium-tuple-space} companion (Section~8) shows the availability
concern dissolves when the space is built \emph{on} a coordinator-free
substrate: the serving node is discovered by capability advertisement, its
failure is detected by evaporation, and a mirror promotes itself through the same
emergent mechanism as every other role in the system. The pull pattern survives;
only the coordinator-shaped infrastructure around it was wrong. (A precision for
readers who know Linda: what survives is generative decoupling and blocking pull,
not associative template retrieval---the realisation in Section~8 is
lane-addressed, with named per-stage queues and opaque payloads, so the claim
path is an O(1) pop rather than a pattern match, and the per-lane depth counters
double as the load signal pull workers route on.)

Carvalho~\cite{carvalho} observes that modern AI agent frameworks are
independently rediscovering Linda's tuple space pattern---implementing ad-hoc
shared stores with task claiming, polling, and matching---without awareness of
forty years of research on its failure modes. The community is not learning from
history. It is repeating it.

\subsection{BDI Agents and FIPA Standards (1990s)}

The Belief-Desire-Intention agent architecture~\cite{bdi} and the Foundation for
Intelligent Physical Agents~\cite{fipa,jade} brought deliberative reasoning to
multi-agent systems. FIPA defined a Directory Facilitator---a registry where
agents register capabilities and query for peers---that is precisely the
registry-based pattern that MIT~NANDA echoes today, four decades later.

\subsection{The Pattern That Persists}

Across five decades, each model correctly identified coordination as a concern
distinct from computation. None eliminated the coordinator assumption:

\begin{center}
\begin{tabular}{ll}
\toprule
\textbf{Model} & \textbf{Coordinator retained} \\
\midrule
Blackboard / HEARSAY (1980) & Central blackboard + control component \\
Actor model (Hewitt, 1973) & Explicit addressing---sender must know recipient \\
Linda (Gelernter, 1985) & Shared tuple space \\
BDI / FIPA / JADE (1990s) & Directory Facilitator as registry \\
Modern LLM orchestration (2020s) & Central orchestrator / CognitiveEngine \\
\bottomrule
\end{tabular}
\end{center}

% ─────────────────────────────────────────────────────────────────────────────
\section{The Coordinator Trap: From Architecture to Symptom}
\label{sec:trap}

We use the term \emph{coordinator trap} to describe the set of failure modes that
any architecture inheriting the coordinator assumption must exhibit. These are not
bugs; they are structural properties. A better coordinator does not escape them.

\subsection{The Mediated Hierarchy Described Precisely}

In a mediated hierarchy, agents are workers in a fanout RPC system. The
coordinator aggregates responses, applies its decision logic, and broadcasts
results. When we call these components ``agents'' we are importing a concept from
the Complex Adaptive Systems literature. Holland~\cite{holland} defines an agent
as an entity that holds a \emph{boundary}---a receptor set specifying the
conditions under which it acts on incoming signals. The boundary is intrinsic to
the agent, not delegated from outside. However, in a mediated hierarchy, agents satisfy
none of this definition: they receive what the coordinator routes, not what they
declare themselves competent to receive. They are passive processors of
coordinator-routed payloads.

\subsection{Audit Burden: The Consequence of Post-Admission Broadcasting}

The unbounded audit obligation follows directly from the coordinator's broadcast
model. Every output of every agent passes through the coordinator and produces a
coordinator-level artifact with identical structure and apparent authority,
regardless of the quality or validity of the underlying agent reasoning. This is
post-admission filtering: artifacts are produced first; quality is assessed after.
As agent populations grow, the coordinator's cognitive load grows linearly with
the number of domains it must evaluate. The audit burden scales with the
coordinator. The three-hour problem and the five-second problem have the same
voice because the coordinator processes them identically~\cite{signal-noise}.

\subsection{Context Loss: The Consequence of Coordinator-Held State}

In a mediated hierarchy, shared state lives in or near the coordinator. Agents are
stateless workers; they hold no durable view of prior decisions. When the
coordinator restarts, the state is gone. When a new agent joins, it must query the
coordinator for a catchup briefing. The attempted fix documented
in~\cite{context-problem}---a filesystem-based shared memory---correctly identifies
that state should be held outside the coordinator, but introduces a shared mutable
store with no eviction policy: entries persist indefinitely, accumulate, and drift
from current system state. This is not a solution to the coordinator trap; it is a
redistribution of the same problem.

\subsection{Output Format Mismatch: The Consequence of Absent Boundary Model}

When the coordinator synthesises a result and broadcasts it, it has no model of
its recipients' boundaries---no representation of which consumers can act on which
signals. All downstream consumers receive the same output because, from the
coordinator's perspective, they are all equivalent consumers of its decisions. A
human reviewer is not equivalent to an agent as a consumer of a coordination
decision, but the coordinator cannot know this without a model of recipient
boundaries.

\subsection{The Category Error: Workers Dressed as Agents}

The three failure modes share a common root: in a mediated hierarchy, components
called ``agents'' are not agents in any theoretically meaningful sense. A component
whose domain is determined by the coordinator's routing decisions is a worker node
in a fanout RPC system. Calling these components ``agents'' inherits none of the
architectural properties that make genuine agents scalable.

The exit from the coordinator trap requires not a better coordinator but the
elimination of the coordinator from the substrate.

% ─────────────────────────────────────────────────────────────────────────────
\section{Theoretical Foundation: Holland's Signal/Boundary Model}

John Holland's framework for Complex Adaptive Systems~\cite{holland,holland-intro},
formalised in \emph{Signals and Boundaries: Building Blocks for Complex Adaptive
Systems} (MIT Press, 2012), provides the theoretical basis for a coordinator-free
architecture.

Holland's thesis: the behaviour of complex adaptive systems emerges from two and
only two primitives.

\textbf{Signals} propagate through a medium unconditionally. No signal is withheld
from propagation on the basis of who might act on it.

\textbf{Boundaries} are receptor sets. Each agent holds a boundary---a set of
conditions under which it \emph{acts} on a signal. Forwarding is always
unconditional; acting is always conditional on boundary admission.

This inversion is the key insight. Conventional distributed systems thinking
routes messages to known recipients; the emitter must know who is listening. In
Holland's model, any agent whose boundary matches a signal responds. The emitter
does not need to know who is listening. The system tolerates churn without stalling
because there is no routing table to maintain, no coordinator to query.

\textbf{Stigmergy and TTL evaporation.} Holland identifies stigmergy---state left
in the medium by agents, readable by other agents---as the mechanism through which
agents coordinate without direct communication~\cite{dorigo}. Applied to
distributed state management, this suggests that shared state should carry a
time-to-live and evaporate unless actively refreshed. An agent that is alive and
relevant keeps its state fresh; a dead agent's state simply fades. Failure is
emergent: the absence of refreshment \emph{is} the failure signal.

Applied to the failure modes of Section~\ref{sec:trap}: pre-admission filtering
structurally prevents the audit burden---an agent whose boundary does not match a
signal produces no artifact to audit. Distributed TTL state makes
coordinator-type context loss structurally highly improbable---there is no
coordinator memory to lose; state evaporation under simultaneous node failure
remains possible but requires active failure of TTL renewal across all replicas
simultaneously. Self-differentiating boundaries make output format mismatch a
design problem rather than a routing obligation. Each failure mode is not
ameliorated; it is either structurally prevented or made structurally highly
improbable at the substrate level.

% ─────────────────────────────────────────────────────────────────────────────
\section{Prior Art: Correct Concepts, Wrong Implementation}

Four prior systems identified the correct underlying concepts but implemented them
with too much protocol ceremony, or in the wrong deployment model, to achieve
their full potential.

\subsection{Jini and the Lease Insight (Sun Microsystems, 1998)}

Sun Microsystems' Java/Jini~\cite{jini-arch,jini-spec} was designed to address
runtime dynamism: in a live network, services appear, depart, and change in ways
no deployment manifest can anticipate, and a crashed service cannot explicitly
deregister. The lease was Jini's answer---decay is the right primitive for an
unpredictably dynamic world. A service holds a lease on its registration; if it
does not renew, the registration expires, providing implicit failure detection
without requiring an explicit deregistration protocol. The insight is correct.
The implementation was protocol-heavy: explicit \texttt{Lease} objects,
\texttt{renew()} RPCs, a lease manager, explicit cancellation. The ceremony
obscures the substrate property---that registrations should evaporate unless
actively maintained---behind an explicit lifecycle protocol.

Jini's other defining mechanism---the dynamic proxy~\cite{jini-arch}---deserves
separate treatment, because its causal direction is the inverse of the one this
paper advocates. A client looking up a service received a service-supplied
object implementing the service interface; inside it, the service chose the
client's transport and serialization. The client's lower layers were payload
shipped from above: the macro specification (the interface) came first, and the
micro realization was manufactured to fit---top-down realization by selection,
in Ellis's sense~\cite{ellis}, where emergence runs the other way. Two
structural consequences followed. First, letting the top configure the bottom
requires a universal executable bottom---a JVM on every participant---so
substrate independence at the interface level was purchased with platform
capture one level down. Second, and more fundamental to this paper's argument:
per-service plumbing fragments the shared medium. The substrate properties that
every emergent layer composes from---unconditional propagation, evaporation,
anti-entropy convergence---are properties \emph{of a common medium}, and a
fleet in which each service ships its own transport has no common medium for
anything to emerge in. Mycelium accordingly keeps Jini's decay insight and
refuses the proxy inversion: the wire format is frozen and identical for every
node, and the only downward influence is stigmergic---macro pheromone fields
constraining locally-deciding agents through the medium, never reconfiguring it.

\subsection{OSGi Requirements and Capabilities}

OSGi's insight was that software agility is a function of modular, dynamically
assembleable components. The OSGi Alliance~\cite{osgi} formalised this as a
dependency model in which software modules declare capabilities they provide and
requirements they need; a resolver matches providers to consumers. The primitive
is correct: declarative matching between providers and consumers, with the
resolver handling wiring.

What mainstream OSGi adoption got wrong was treating resolution as
static---performed once at bundle-install time. A module that disappeared at
runtime left a gap with no in-process mechanism for repair; closing it required
a human to rebuild and redeploy. Dynamic assembly was the goal; deploy-time
resolution was where that ambition stopped. Paremus Service Fabric addressed
this directly.

\subsection{Paremus Service Fabric and the Reconciliation Engine}

Paremus Service Fabric (circa 2010--2015)~\cite{paremus} began as a Jini runtime
with OSGi components, using Jini's service discovery as the underlying discovery
mechanism. Over time, Jini service discovery was replaced with the OSGi Remote
Services specification backed by an in-house gossip-based
implementation---decentralised discovery, TTL-managed registrations, no central
lookup service. What emerged from this evolution was a deliberate fusion of
Jini's runtime dynamism with OSGi's modular assembly model: the R\&C graph as a
\emph{declared target topology}, with resolution made continuous against a live,
gossip-discovered environment rather than a static deployment manifest. The
runtime was continuously monitored against that declared target state and deltas
driven back into convergence---closer to the Kubernetes control loop than to a
shared knowledge graph, and requiring no human redeploy.

What did not change was the hermetic relationship between the gossip layer and
the reconciliation engine: `as above, so below'---the declared target state
lived in the central engine above; the gossip layer revealed the landscape below;
the engine held the authority to reconcile the two. The coordinator trap persisted
not in the discovery mechanism but in the architectural layer that held the
target.

What Paremus still required was a \emph{central reconciliation engine} holding
that target state. If it went down, reconciliation stopped. The coordinator trap
in a more sophisticated form.

\subsection{Flow-Based Programming (Morrison, c.~1971)}

Flow-Based Programming~\cite{fbp}, developed by J.~Paul Morrison at IBM in the
early 1970s, conceived programs as networks of independent black-box processes
communicating through bounded connections. The model separates coordination from
computation and gets three things right: (i)~\emph{reusable processes} that know
nothing about their neighbours; (ii)~\emph{automatic backpressure}, in which
bounded connections block writers when full; and (iii)~\emph{declarative topology}
expressed as a data structure separate from component code.

What FBP retained was an external medium to which all these properties belong. The
FBP runtime \emph{is} the coordinator: it allocates connections, manages
backpressure, and resolves the wiring. The coordinator trap remains.

\subsection{The Strip-the-Ceremony Pattern}

A pattern emerges across all four cases:

\begin{center}
\renewcommand{\arraystretch}{1.25}
\begin{tabularx}{\textwidth}{>{\bfseries}lXXX}
\toprule
Prior art & Correct concept & Implementation ceremony & Substrate property equivalent \\
\midrule
Jini & Registrations should decay & \texttt{Lease.renew()} RPC + lease manager & TTL as natural evaporation \\
OSGi & Declarative capability matching & Static bundle-install resolver & Continuous evaluation against live state \\
Paremus & Continuous reconciliation toward target state & Central reconciliation engine + target state graph & Gossip mesh --- every TTL refresh is a reconciliation tick \\
FBP & Backpressure, declarative topology, reusable processes & External FBP runtime managing connections & Opacity KV flag --- agent self-declares saturation; resolvers skip it \\
\bottomrule
\end{tabularx}
\end{center}

The pattern that produces better architecture: identify the correct concept, find
the substrate \emph{property} that produces the same behaviour without an explicit
protocol, implement the property and let the behaviour emerge.

% ─────────────────────────────────────────────────────────────────────────────
\section{Mycelium: A Coordinator-Free Substrate}
\label{sec:mycelium}

Mycelium is an embeddable Rust library implementing Holland's signal/boundary
model. It has no daemon, no control plane, no installer, no orchestrator. It is a
crate embedded in the caller's process. The architecture is three layers; each directly addresses one of the three
structural failure modes identified in Section~\ref{sec:trap}: audit burden
(\S4.2), context loss (\S4.3), and output format mismatch (\S4.4).

\subsection{Layer~I --- KV Store: Addressing Context Loss}

Layer~I is a gossip-replicated key-value store. Every entry carries a Hybrid
Logical Clock timestamp for causal last-write-wins ordering and a time-to-live.
Entries are replicated across the cluster via epidemic gossip; anti-entropy
reconciliation runs on peer reconnection. There is no primary, no leader, no
coordinator. Every node holds a full replica of the key namespace relevant to it.

\textbf{Context loss structurally highly improbable.} Context loss in the mediated
hierarchy arises because shared state is held in or near the coordinator. In
Mycelium's Layer~I, there is no coordinator memory to lose. A node that restarts
reconnects to its peers, runs anti-entropy, and recovers full mesh state within
one gossip cycle. State loss requires simultaneous TTL expiry across all
replicas---possible under full-cluster failure, but structurally different from
the single-point loss inherent in coordinator designs.

\textbf{TTL evaporation as stigmergy.} Every capability advertisement, every group
membership record, every opacity flag is a KV entry with a TTL. A live node
continuously re-advertises; the trail stays fresh. A dead node stops
re-advertising; its entries evaporate. There is no explicit failure detection
protocol because absence of refreshment \emph{is} the failure signal.

\textbf{The key inversion from Paremus.} In Paremus, the target state was held
\emph{above} the runtime in a central graph and pushed down into components. In
Mycelium, the target state is \emph{compiled into the runtime components
themselves}---each node carries its own fragment as capability and requirement
declarations. There is nothing external to converge toward. Convergence emerges
upward from those fragments.

\subsection{Layer~II --- Signal Mesh: Addressing Audit Burden and Output Mismatch}

Layer~II is an ephemeral signal mesh implementing Holland's signal/boundary model
directly. Signals are emitted with a scope and propagated unconditionally to all
reachable nodes. At each receiving node, \texttt{Boundary::admits()} evaluates the
signal against the node's declared receptor set. Only nodes whose boundary matches
the signal act on it. Forwarding is always unconditional; acting is always
conditional on boundary admission.

\textbf{Audit burden structurally prevented.} A signal scoped to the
\texttt{gpu-compute} group only activates nodes in that group. Signals propagate
at the transport layer unconditionally---that is the substrate's routing
mechanism---but an agent in an unrelated group produces no \emph{application-layer
artifact}: it does not act, does not log, does not emit a downstream result. The
audit burden that scales with a coordinator scales because every event passes
through a single processing point that generates a record regardless of
relevance. Here, the record is never created because the admission test fails
before processing begins. The three-hour problem and the five-second problem
never land on a reviewer who cannot evaluate them.

\textbf{Output format mismatch structurally addressed.} A signal emitted to a
\texttt{human-review} group reaches only nodes whose boundary includes
human-review membership; a signal emitted to \texttt{automated-pipeline} reaches a
different set. Human-facing and agent-facing consumers self-differentiate via their
boundary declarations without any routing logic in the emitter. The
representation design obligation---what a human-review signal should
\emph{contain}---remains with the emitting agent; the substrate ensures it
reaches the right consumers, not that its content is appropriate for them.

\textbf{Causal ordering.} For applications requiring strict per-sender ordering,
\texttt{emit\_ordered()} extends the signal mesh with a receiver-side causal
reorder buffer. Each ordered message carries an HLC sequence counter in the wire
frame (introduced at wire protocol version~11); the receiver reassembles messages
per sender in HLC order before delivering them to handlers. Applications tolerating
reorder use the standard \texttt{emit()}; those requiring strict causal delivery
opt in via \texttt{emit\_ordered()}. The common case---ephemeral broadcast
signals---pays no overhead, because ordering is applied only when explicitly requested.

\subsection{Capability and Requirement Declarations: Addressing the Category Error}

Each Mycelium node holds a \texttt{CapabilityGroupDef}---a declarative
specification of capabilities it provides and requirements it needs. The node
independently evaluates whether it should join each capability group by testing its
own capabilities against the group's filter, right now, against live KV state.
Membership is self-assigned, not coordinator-assigned. The node \emph{is} its
boundary. The application's coordination topology emerges from the composition of
these declarations across the live node population, not from a coordinator's
routing table.

\subsection{Layer~III --- Consensus: Addressing the Coordinator Bottleneck}

Layer~III provides epidemic consensus---\texttt{group\_propose},
\texttt{system\_propose}, and \texttt{cross\_group\_propose}---implemented entirely
through the signal mesh. \texttt{PROPOSE}, \texttt{VOTE}, and \texttt{COMMIT} are
signal payloads riding ordinary Layer~II frames. The substrate is unaware that
consensus is happening.

\textbf{Fault model and guarantees.} The consensus layer is crash-fault tolerant;
it does not provide Byzantine fault tolerance. A ballot commits when a configurable
quorum fraction of group members---defaulting to a simple majority---has voted. A
node that crashes mid-ballot simply stops voting; the ballot either reaches quorum
without it or times out. Under network partition, liveness is not guaranteed: a
minority partition will not commit. Safety---no two nodes commit conflicting values
for the same ballot---holds as long as quorum is not simultaneously satisfied by
two disjoint subsets, which a majority quorum prevents. Quorum membership is
evaluated against the live capability-group population at ballot time, not a fixed
initial roster; nodes that join or leave between proposal and commit may affect
whether quorum is reachable, but cannot cause two conflicting values to commit.
This is not a replacement for Raft, Paxos, or BFT consensus in high-assurance
settings; it is a coordination primitive for capability-group decisions in
homogeneous trusted clusters where crash-fault tolerance and leaderless operation
are the primary requirements.

\texttt{cross\_group\_propose} supports ballots where multiple named capability
groups act as independent voting blocs. A proposal commits only when every group
independently reaches its required quorum fraction. This supports multi-AZ
durability requirements, compliance ratification, and hierarchical AI pipelines---all
without a coordinator.

\textbf{Coordinator bottleneck structurally addressed.} In a mediated hierarchy,
coordination throughput is bounded by the coordinator's capacity, and coordinator
failure stops coordination entirely. In Layer~III, there is no coordinator. No
single node's failure prevents other ballots from proceeding. The system degrades
gracefully to reduced quorum, not coordination failure.

\subsection{Agentic Flow Networks}

The three-layer architecture described in \S\S7.1--7.4 supports a distinct
multi-agent computational pattern that we term an \emph{Agentic Flow Network}
(AFN). An AFN is a topology of capability-bounded agent groups in which work flows
through the named lanes of the tuple-space companion (Section~8), each stage
\texttt{take}-ing from one lane and \texttt{complete}-ing atomically into the
next. The pipeline is data topology, not node topology: a worker is ``at'' a
stage only for the duration of one take~$\to$~process~$\to$~complete cycle.

An AFN has five structural properties:

\begin{enumerate}
\item \textbf{Substrate Unity.} All \emph{coordination}---capability
  advertisement, serving-role election, group membership, backpressure pheromones,
  completion markers---is carried by the same gossip KV substrate; the work-item
  payload path is a point-to-point RPC to the space's serving node (single-copy,
  O(1) per item, versus O(N) gossip replication). There is no separate message
  queue, registry, or coordination bus.

\item \textbf{Topology Emergence.} Which nodes participate is not configured by
  an operator. Workers self-select stages by reading per-lane depth; the space's
  serving roles---primary and mirroring secondary---are discovered by capability
  advertisement and re-elected through evaporation when a holder fails.

\item \textbf{Fluid Allocation.} There is no static assignment of workers to
  stages. Each worker's flow loop takes from the deepest lane in its repertoire,
  so the pool masses on the bottleneck automatically---fluidity is self-selection
  against pressure.

\item \textbf{Depth-Native Backpressure.} Per-lane depth counters are
  simultaneously the backpressure signal (\texttt{put} is refused beyond a high
  watermark) and the allocation signal workers route on. Saturation is observable
  state in the medium, not a flow-control protocol.

\item \textbf{Deadline-Native Recovery.} Items taken but not acknowledged
  re-queue automatically when the in-flight deadline lapses, and a stage
  transition is a single WAL record, so a crash cannot half-replay a hop. Absence
  of acknowledgement is the failure signal, exactly as absence of refreshment is
  for capabilities.
\end{enumerate}

These properties are validated empirically by \texttt{examples/fluid\_pipeline/}
---a 10-worker, 4-stage news article pipeline (parse $\to$ enrich $\to$ score
$\to$ aggregate) whose workers take from the deepest stage and complete into the
next, exercised end-to-end in CI in both its pull form and its retained baseline,
alongside integration scenario~13 (the cross-node put/take/complete/ack
lifecycle) and scenario~11 (the substrate-level stage-transition properties). An
earlier formulation used the replicated KV ring itself as the inter-stage
buffer; it survives in the example as the push-mode comparison baseline that the
three-arm experiment (Future Work) requires.

An AFN is not a new programming model layered on top of Mycelium. It is the name
for the pattern that naturally arises when a distributed application is built
correctly on a substrate that eliminates the coordinator.

% ─────────────────────────────────────────────────────────────────────────────
\section{Evaluation}

Mycelium's correctness across its three-layer architecture is validated by 305
unit tests and 13 integration scenarios run against a live 5-node Docker cluster.
Scenarios cover KV replication under partition and reconnection, signal delivery
and boundary admission, capability group formation and dissolution, consensus
quorum under node failure, cross-group voting, the full Agentic Flow Networks
pipeline, Prompt Skills cross-node KV propagation with LLM invocation, and the
tuple-space pull pipeline described below. All 13 scenarios pass at HEAD.

\textbf{The tuple-space companion as constructive evidence.} The claim of
Section~4.5---that the exit from the coordinator trap is elimination, not
improvement---admits a constructive test: rebuild the most coordinator-shaped
workload, work distribution, with no component that predicts any agent's state.
The \texttt{mycelium-tuple-space} companion crate does this using only the
substrate's public API (the dependency relation itself demonstrating that pull
composes from the existing primitives). Producers \texttt{put} items into named
stages---per-stage FIFO lanes, not Linda's associatively-matched bag; payloads
are opaque to the space and an item's pipeline position is the lane it
occupies. Workers block on \texttt{take} and receive work exactly when ready---
readiness is announced by the claim itself, so the staleness that drives
misrouting in push systems has no variable to inhabit. Three test results carry
the weight: (i)~1{,}000 items across ten concurrent producers and ten concurrent
workers are delivered exactly once, with no participant ever holding a view of
another's load; (ii)~when the serving node is killed with unacknowledged items
live, a mirroring standby detects the capability evaporation, promotes itself
through the same emergent election as any other role, and serves the surviving
items under their original identifiers---acknowledged items do not resurrect;
(iii)~a full cross-node integration scenario drives the
put/take/complete/ack lifecycle through the HTTP gateway, including in-flight
accounting and the blocking-take timeout contract. The historical objection to
tuple spaces---the shared medium as a designed-in coordinator requiring its own
consistency machinery (Section~3.3)---is answered by construction: the rendezvous
point's availability is maintained by substrate evaporation and emergent
promotion, and its decision authority is zero because it makes no decisions.

The implementation is publicly available at
\url{https://github.com/RichardEko/mycelium} (tag \texttt{paper-submission-v2}) under the AGPL-3.0 licence.
The integration test harness is included in the repository and reproducible with a
single \texttt{make test} invocation against a Docker-composed 5-node cluster.
The AGPL-3.0 licence is intentional: applications that deploy Mycelium as part of a
networked service must publish their source under the same terms, closing the
SaaS loophole that permissive licences leave open. Organisations requiring
proprietary embedding may obtain a commercial licence from Tathata Systems Ltd.

% ─────────────────────────────────────────────────────────────────────────────
\section{Discussion}

\subsection{The Hayek Parallel}

The coordinator trap is not a new discovery. Friedrich Hayek described it in
1945---for economies.

In \emph{The Use of Knowledge in Society}~\cite{hayek}, Hayek argued that the
central planning problem is not computational---it is epistemic. No central
planner can possess the distributed, local, tacit knowledge held by individual
market participants. Prices are not just numbers; they are signals that aggregate
and propagate dispersed knowledge through the economy without anyone needing to
understand the whole.

The parallel to the coordinator trap is suggestive rather than exact---economies,
biological systems, and distributed software are not identical. But all three
expose the same cost of centralising knowledge that is generated and updated at
the edge. A planned economy and a mediated hierarchy share a recognisable failure
mode: a central node that must aggregate knowledge it cannot fully possess,
synthesise decisions on behalf of participants who understand their local context
better than any coordinator can, and issue commands downward into a system whose
ground truth is always more current at the edges than at the centre.

Hayek's market is a signal/boundary system. He just did not have Holland's
vocabulary.

\subsection{The Beinhocker Parallel: Organisations as Complex Adaptive Systems}

Beinhocker's \emph{The Origin of Wealth}~\cite{beinhocker} reaches the same
conclusion from a third direction: organisations. Drawing on extensive empirical
analysis of corporate performance, Beinhocker finds that hierarchical organisations
centralising strategic information and decision-making systematically underperform
distributed ones---not because their leaders are less capable, but because the
local knowledge required for adaptive decisions cannot travel up the chain fast
enough to remain actionable by the time it arrives.

Hayek, Holland, and Beinhocker are describing the same property of complex
adaptive systems at three different scales---economies, biological systems, and
organisations. The coordinator trap is not a software engineering peculiarity. It
is a structural consequence of centralising coordination in any sufficiently
complex adaptive system.

\subsection{The Economic Case: Quadratic Cost Decomposition}

The architectural argument above is independent of cost. There is also a direct
economic argument. For workloads in which context accumulates across
iterations---the common case
in long-running coordinator pipelines where each step appends to a growing
history---LLM agent execution cost scales superlinearly with iteration count.
Because each API call re-processes the entire accumulated context, the token
count at iteration~$k$ is $\Delta T + \tau k$ (where $\Delta T$ is the initial
prompt and $\tau$ is context growth per iteration). Summing over $k$ iterations
yields a total proportional to $k^2$: cost is quadratic in iteration count, not
linear. Systems that summarise, truncate, or selectively retrieve context can
reduce but not eliminate this accumulation; the model below applies at its
strongest to pipelines that maintain a running context window.

The decomposition leverage is direct. Splitting a $k$-iteration task into $M$
independent subtasks of $k/M$ iterations each reduces total cost by factor~$M$:
\[
  M \times \left(\frac{k}{M}\right)^2 = \frac{k^2}{M}
\]
versus $k^2$ for the monolithic run. A 4-way decomposition of a 100-iteration task
costs $4 \times 25^2 = 2{,}500$ versus $100^2 = 10{,}000$---a 4$\times$ saving
from a 4$\times$ decomposition.

A mediated hierarchy cannot realise this saving without explicit orchestration
overhead. The coordinator's synthesis step, where it aggregates $M$ subtask
results into a combined output, recreates the full accumulated context and pays
the quadratic penalty once more. However, Mycelium's capability-bounded groups
implement decomposition automatically: each group's TTL constrains how long a work
item remains live in that stage, which under typical workloads approximates a bound
on effective $k$ for that group---a time bound rather than a strict iteration count
bound, but one that produces the same practical effect when per-iteration latency
is roughly constant. The quadratic cost accumulation resets at each group boundary,
and no coordinator synthesis is required.

The practical significance of this saving is reinforced by empirical energy
measurements. White~\cite{white2025} demonstrates that inference energy per
query varies by two orders of magnitude across model families and sizes, and
that energy cost is highly sensitive to input characteristics. Where the
quadratic accumulation compounds an already expensive per-call baseline, the
saving from decomposition is not a theoretical nicety---it is an operational
necessity.

\subsection{Limitations}

Mycelium assumes a cluster the operator owns. Mycelium's scope is intra-cluster;
cross-organisational discovery requires external registry infrastructure.
Mycelium's conforming A2A endpoint makes it reachable from any A2A-speaking
registry without modification to either side.

Ephemeral signals are intentionally not durable---a node that misses a signal
misses it. Durable delivery is a higher-order concern built on the KV layer or
consensus, not a substrate property.

Boundary admission requires agents to declare their boundaries correctly. A
misconfigured boundary---too broad or too narrow---produces incorrect routing
without any coordinator to catch the error. This places a correctness obligation
on the capability declarations that the mediated hierarchy places on the
coordinator instead.

\textbf{Trust and adversarial behaviour.} Self-declared capability boundaries
raise questions about false capability claims, malicious group admission, and
poisoned KV entries. Mycelium's intra-cluster assumption---the operator owns the
cluster---substantially reduces this surface, but owned clusters still contain
buggy or compromised nodes. The \texttt{tls} feature addresses node identity:
each node holds an Ed25519 keypair; mTLS authenticates every TCP connection;
KV writes under the \texttt{tls} feature are Ed25519-signed at the frame level
(introduced at wire protocol version~10; current protocol version is~11),
making undetected mutation of replicated state infeasible for
an attacker who has not compromised the originating node's private key. The
current implementation does not provide capability-claim attestation---a node
can declare capabilities it does not possess---and Byzantine-tolerant consensus
is out of scope. These are correct tradeoffs for the target deployment model
(trusted, operator-controlled clusters) but should be revisited before Mycelium
is used in adversarial multi-tenant environments.

The quadratic cost decomposition in \S9.3 assumes that the $k$-iteration task
can be partitioned into $M$ independent subtasks. Many real agentic workloads
are inherently sequential---multi-step reasoning chains, document drafting,
iterative refinement---where each step depends causally on the prior and
decomposition is not valid. The saving applies where subtask independence holds;
it does not apply universally.

\textbf{When not to use AFN.} Five cases warrant a different approach:
\emph{Strict ordering}---where tasks must be processed in a guaranteed sequence,
not by the first available worker.
\emph{Exactly-once delivery}---where each work item must be processed exactly
once, with no duplicate execution under any failure scenario.
\emph{Complex DAG dependencies}---fan-in joins (a stage consuming the join of
two upstream outputs) are not expressible in exactly-named lanes today; the
keyed-exact-match \texttt{take} discussed below (\texttt{take\_by\_key}, shipped as
M13 in v2.0, 2026-06-20) covers pairwise joins, and arbitrary cross-stage predicates belong to the blackboard
companion (Future Work). An orchestrated DAG engine (Airflow, Dagster) remains
right where centralised control is itself the requirement---human approval
gates, one pane for dependency state---a coordinator adopted deliberately.
\emph{Long-term log retention}---where TTL evaporation is a correctness property;
applications requiring replay capability must use a separate append-only store.
\emph{Cross-cluster fan-out at scale}---where the gossip mesh operates within a
single cluster; multi-cluster deployments require a federation layer.

\textbf{Lane-addressed claiming, and where Linda's matching earns its keep.}
The Section~8 realisation replaces associative template retrieval with named
per-stage lanes, and the honest statement of that choice is a decomposition
principle rather than a deletion: the architecture retains associative matching
at the \emph{control plane}---the capability resolver's attribute filters and
the signal boundary's receptor predicates are both template matchers,
re-evaluated continuously---and keeps it off the \emph{durable claim path},
where a matching engine would un-buy the lane properties (O(1) claims,
well-defined per-lane depth and therefore backpressure, and the single-record
stage transition). Most workloads decompose along that line: match to select a
lane, pull from the lane. Four do not, and they mark the realisation's boundary
of validity: (i)~\emph{fan-in joins} degenerate lanes to one item per
correlation key---the keyed-exact-match \texttt{take\_by\_key} covers this without
general matching; (ii)~\emph{contract-net claiming}, where a multi-dimensional
predicate lives on the item, makes per-item matching irreducible;
(iii)~\emph{blackboard working memory}---opportunistic reasoning over shared
typed facts, with no stages because the flow topology is emergent per item; the
one missing primitive is a competitive destructive claim-by-predicate with the
tuple space's exactly-once discipline, and this is precisely the pattern
Carvalho~\cite{carvalho} observes agent frameworks rebuilding ad hoc;
(iv)~\emph{semantic claiming}---similarity between task content and agent
skill profile---is a ranking concern for the selection edge, never the claim
path. Consistent with the layering principle, the answer to all four is
composition above the substrate, not a matching engine inside it.

\subsection{Future Work}

\begin{itemize}
\item \textbf{Empirical comparison} against a deployed mediated hierarchy at
  equivalent agent counts. Key measurements: (i)~coordination convergence
  time---single-ballot \texttt{group\_propose} vs NegMAS~\cite{negmas} SAO
  N-round negotiation; (ii)~failure tolerance---coordinator failure in a mediated
  hierarchy vs random node failure in Mycelium; (iii)~state freshness under
  churn---TTL evaporation latency vs knowledge-graph drift rate; (iv)~audit
  obligation under load---artifact count growth as agent population scales.

\item \textbf{Completing the three-arm work-distribution experiment.} Two of
  the three arms have already been run head-to-head on this substrate:
  broker-mediated prediction versus locally evaluated prediction, with
  decision-level results (latency, staleness, misroute rate at
  $N=10$--$40$) reported in the companion paper~\cite{companion} and the harness in the
  repository. The third arm---pull---is implemented as the tuple-space
  companion (Section~8) but has not yet been compared against the other two,
  because it \emph{cannot} be compared on the same instruments: pull abolishes
  the staleness/misroute vocabulary, so the unified experiment must measure
  outcomes---end-to-end job latency, throughput, idle-while-work-exists time,
  and fairness---under heterogeneous and drifting worker service times. The
  companion paper's framework predicts the prediction arms degrade with
  heterogeneity while pull does not; that prediction is what remains untested.

\item \textbf{Formal verification} of the signal/boundary substrate properties
  using TLA+ or similar.

\item \textbf{Cross-cluster federation}---practical experience deploying multiple
  independent Mycelium clusters registered with an internet-scale A2A
  gateway~\cite{nanda}, measuring discovery latency and trust propagation across
  organisational boundaries.

\item \textbf{Byzantine fault tolerance}---the current consensus layer is
  crash-fault tolerant only. Extending it to tolerate actively malicious nodes
  while preserving the leaderless, coordinator-free property is an open problem.

\item \textbf{A blackboard companion as a second constructive test.} The
  tuple space proves the no-coordinator claim for \emph{pipeline} work, where
  every item's next step is known in advance. The natural second test is
  \emph{opportunistic} work: a shared pool of facts---the blackboard of
  Section~3.1, reborn today as LLM-agent shared scratchpads---where any agent
  whose trigger condition matches a posted fact reacts to it, and the flow
  topology emerges per item. The substrate already provides most of this
  (facts propagate to everyone; triggers are each agent's own local
  declarations); the single missing piece is an atomic \emph{claim} of a
  matching fact, so that when two agents' triggers both fire, exactly one
  consumes it. (Concretely: in a community-microgrid scratchpad, any number of
  agents may freely \emph{read} a posted surplus-solar fact and update their
  forecasts and prices---but when a ``store 3\,kWh from feeder~4'' task
  appears and two battery agents both match it, exactly one must win it: the
  energy exists once. Reading is concurrent; acting is exactly-once.) (Concretely: in an incident-response scratchpad, any number of
  diagnostic agents may freely \emph{read} a posted alert and add
  observations---but when a ``roll back the deploy'' task appears and two
  executor agents both match it, exactly one must win it. Reading is
  concurrent; acting is exactly-once.) Building that as a companion crate---on the public API, with
  the tuple space's durability discipline---would extend the Section~8
  constructive proof to the workload class at the lane model's boundary, and
  to exactly the pattern agent frameworks are rebuilding ad
  hoc~\cite{carvalho}. This is design-sketch stage in the repository, not yet
  implemented.
\end{itemize}

% ─────────────────────────────────────────────────────────────────────────────
\section{Conclusion}

The coordinator is not a solution to the multi-agent coordination problem. It is a
restatement of the problem at a different scale. Mediated hierarchies make the
coordinator responsible for filtering, memory, routing differentiation, and fault
tolerance---properties that in a correctly designed substrate belong to the agents
themselves.

Unbounded audit obligation, context loss across restarts, and output format
mismatch between heterogeneous consumers are not bugs in specific implementations.
They are structural consequences of the coordinator assumption, predictable from
first principles, and irreducible by improving the coordinator.

Holland's signal/boundary model provides the theoretical basis for a different
architecture: one in which coordination emerges from substrate properties rather
than explicit protocols. Mycelium implements this model as an embeddable library.
Each of its three layers---gossip KV store, signal mesh with boundary admission,
epidemic consensus---addresses a mirrored failure mode, not by handling the
failure gracefully but by making each failure mode structurally prevented or,
where substrate constraints apply, structurally highly improbable.

The target state for a Mycelium application is not a document held in a registry
or a graph maintained by a reconciliation engine. It is the aggregate of what
every node declares itself to be, compiled into the runtime components at build
time, assembled bottom-up by the mesh at runtime, and always current because
anything not actively maintained evaporates. Coordination emerges. No coordinator
required.

The coordinator trap is not a new discovery. Hayek described the epistemic cost
of centralisation for economies in 1945. Holland formalised the signal/boundary
model for complex systems in 2012. Mycelium demonstrates that both insights
transfer to AI agent fleets. The claim is the same in all three cases:
distributed local knowledge, expressed through signals and boundaries, produces
emergent order that no single coordination point can replicate---because the
knowledge required is always more current at the edges than at the centre.

% ─────────────────────────────────────────────────────────────────────────────
\bibliography{references}

% ─────────────────────────────────────────────────────────────────────────────
\section*{About the Author}

\textbf{Dr.\ Richard Nicholson} is CEO and founder of Tathata Systems Ltd and
Chief AI Transformation Officer at Novus-i2, a strategic transformation company
based in the United Kingdom. He was the founder and CEO of Paremus Ltd., the
company behind Paremus Service Fabric---the continuous OSGi Requirements \&
Capabilities runtime resolver discussed in Section~6.3 as direct conceptual
predecessor to Mycelium's capability model. First-hand experience designing and
deploying Service Fabric over a decade of production use informed both the
architectural critique presented in this paper and the design decisions made to
move beyond it.

He served as President, Board Director, and Treasurer of the OSGi Alliance for
approximately six years~\cite{osgi-mod,osgi-cmb}. He was awarded the OSGi Laureate
(2019) for outstanding contributions to the Alliance and the OSGi Leadership Award
(2019) for service as a board officer and advocate for OSGi modularity strategy.
He was co-architectural lead for the EU Horizon~2020 Brain-IoT project, a
federated platform for emergent coordination across heterogeneous IoT and edge
devices~\cite{brain-iot-1,brain-iot-2}.

\end{document}
